\section{Experiments}
We empirically demonstrate that DBB estimation is effective in the RLVR setting. Section~\ref{subsec:expsetting} describes the models, datasets, baselines, and detailed setup used in our experiments, while Section~\ref{subsec:mainresult} presents the experimental results.

\subsection{Experimental Settings}
\label{subsec:expsetting}

\paragraph{Models \& Datasets.}
We conduct experiments on two model sizes, Qwen3-1.7B-Base and Qwen3-8B-Base \cite{yang2025qwen3technicalreport}, to investigate the effect of model scales.
For both model scales, we use DAPO-Math-17k \cite{yu2025dapoopensourcellmreinforcement} as the training dataset.
As in-distribution evaluation benchmarks for mathematical reasoning, we consider six widely used datasets: MATH500 \cite{math500}, Minerva \cite{minerva}, AIME24/25, AMC24 \cite{li2024numinamath}, and OlympiadBench (Olympiad; \citeauthor{olympiad}, \citeyear{olympiad}).
To further assess out-of-distribution generalization, we evaluate performance on MMLU-Pro \cite{mmlupro}, GPQA-Diamond (GPQA-D; \citeauthor{gpqa}, \citeyear{gpqa}), and Big-Bench Hard (BBH; \citeauthor{bbh}, \citeyear{bbh}).

\paragraph{Baselines.}
The DBB reward estimation aims to improve reward distribution estimation by partially leveraging historical information.
To isolate the effect of DBB, we first compare GRPO with the DBB estimator (GRPO-DBB) and naive GRPO, which relies on point estimation.
In addition, we include three representative variance-collapse-aware or history-based baselines: RePO~\citep{li2025reporeplayenhancedpolicyoptimization}, which replays full historical trajectories; GRESO~\citep{zheng2025actpaysefficientreinforcement}, which probabilistically filters prompts likely to lead to variance collapse; and DAPO~\citep{yu2025dapoopensourcellmreinforcement}, which uses dynamic sampling to discard zero-variance prompts.
For all baselines, we follow the best configurations reported in their respective papers.
To ensure a fair comparison, we standardize the total number of rollouts used throughout training across all methods.



\paragraph{Training \& Evaluation Setups.}
All training experiments are conducted using the \texttt{verl}\footnote{\url{https://github.com/verl-project/verl}} framework, and \texttt{Math-Verify}\footnote{\url{https://github.com/huggingface/Math-Verify}} is employed to extract and verify final answers.
Experiments with Qwen3-1.7B-Base and Qwen3-8B-Base are conducted on 4$\times$H200 and 8$\times$H200 GPUs, respectively, except for RePO.
Due to the requirement for additional GPU memory, RePO is trained on 8$\times$H200 GPUs for both model scales.
The rollout batch size is set to 128, and the gradient update batch size is set to 64.
We sample 8 responses per query from the on-policy model and train the model for a total of four epochs.
For evaluation, we conduct four independent runs per method with different random seeds, sampling 8 responses per query in each run (i.e., 32 responses per prompt in total), and report Acc@8 as the evaluation metric.
Each cell reports the mean and standard deviation across the four runs, and statistical significance is assessed via one-sided paired $t$-tests over per-prompt accuracies.
The training setup details for each baseline and the sampling parameters used for validation and evaluation are also provided in Appendix~\ref{ap:A}.



\subsection{Main Results}
\label{subsec:mainresult}
\begin{table*}[t]
\centering
\scriptsize
\setlength{\tabcolsep}{2pt}
\caption{
In-distribution (ID) and out-of-distribution (OOD) evaluation results trained on DAPO-Math-17k.
For a fair comparison, all methods are trained with the same total number of rollouts.
Evaluation is conducted over four independent runs with Acc@8 per run (i.e., 32 responses per prompt in total); each cell reports mean $\pm$ standard deviation across the four runs.
$\Delta$ rows report the absolute Acc@8 improvement of GRPO-DBB over each baseline together with a one-sided paired $t$-test $p$-value ($H_1$: GRPO-DBB $>$ baseline).
}
\label{tab:main_results}
\resizebox{\textwidth}{!}{
    \begin{tabular}{lcccccc>{\columncolor{AvgCol}}c|ccc>{\columncolor{AvgCol}}c}
    \toprule
    & \multicolumn{7}{c|}{\textbf{In-Distribution}}
    & \multicolumn{4}{c}{\textbf{Out-of-Distribution}} \\
    \cmidrule(lr){2-8} \cmidrule(lr){9-12}
    Method
    & MATH500 & Minerva & AIME24 & AIME25 & AMC24 & Olympiad & Avg.
    & MMLU-Pro & GPQA-D & BBH & Avg. \\
    \midrule

    \rowcolor{ModelRow}
    \multicolumn{12}{c}{\textbf{Qwen3-8B-Base} trained with \textbf{DAPO-Math-17k}} \\
    \midrule
    GRPO
    & 87.87{\tiny$\pm$0.36} & 39.34{\tiny$\pm$0.29} & 29.17{\tiny$\pm$0.68} & 25.31{\tiny$\pm$2.16} & 57.71{\tiny$\pm$1.44} & 54.54{\tiny$\pm$0.19} & 48.99{\tiny$\pm$0.46}
    & 50.16{\tiny$\pm$0.33} & 45.27{\tiny$\pm$1.69} & 59.18{\tiny$\pm$0.33} & 51.54{\tiny$\pm$0.55} \\
    RePO
    & 86.22{\tiny$\pm$0.31} & 35.02{\tiny$\pm$0.47} & 27.29{\tiny$\pm$0.80} & 23.44{\tiny$\pm$1.04} & 56.25{\tiny$\pm$0.70} & 52.70{\tiny$\pm$0.44} & 46.82{\tiny$\pm$0.07}
    & 53.29{\tiny$\pm$0.55} & 46.34{\tiny$\pm$1.74} & 42.44{\tiny$\pm$0.30} & 47.35{\tiny$\pm$0.80} \\
    GRESO
    & 87.01{\tiny$\pm$0.20} & 33.88{\tiny$\pm$0.60} & 27.36{\tiny$\pm$0.79} & 22.78{\tiny$\pm$1.09} & 57.31{\tiny$\pm$0.57} & 52.81{\tiny$\pm$0.17} & 46.86{\tiny$\pm$0.12}
    & 62.81{\tiny$\pm$0.39} & 48.53{\tiny$\pm$0.43} & 65.17{\tiny$\pm$0.63} & 58.83{\tiny$\pm$0.40} \\
    DAPO
    & 86.00{\tiny$\pm$0.07} & 37.08{\tiny$\pm$0.30} & 31.06{\tiny$\pm$1.19} & 21.02{\tiny$\pm$1.14} & 52.67{\tiny$\pm$0.92} & 51.64{\tiny$\pm$0.14} & 46.58{\tiny$\pm$0.16}
    & 59.28{\tiny$\pm$0.67} & 45.05{\tiny$\pm$1.79} & 63.41{\tiny$\pm$0.29} & 55.91{\tiny$\pm$0.46} \\
    \midrule
    GRPO-DBB
    & \textbf{88.71}{\tiny$\pm$0.18} & \textbf{39.46}{\tiny$\pm$0.51} & \textbf{33.60}{\tiny$\pm$2.05} & \textbf{29.02}{\tiny$\pm$2.53} & \textbf{60.54}{\tiny$\pm$1.00} & \textbf{56.53}{\tiny$\pm$0.25} & \textbf{51.31}{\tiny$\pm$0.53}
    & \textbf{63.12}{\tiny$\pm$0.38} & \textbf{49.43}{\tiny$\pm$2.11} & \textbf{67.09}{\tiny$\pm$0.13} & \textbf{59.88}{\tiny$\pm$0.66} \\
    $\Delta$ w.r.t.\ GRPO
    & +0.84{\tiny\,(0.0073)} & +0.12{\tiny\,(0.2015)} & +4.44{\tiny\,(0.0114)} & +3.71{\tiny\,(0.0063)} & +2.83{\tiny\,(0.0018)} & +1.99{\tiny\,(0.0009)} & +2.32{\tiny\,(0.0007)}
    & +12.96{\tiny\,(0.0000)} & +4.17{\tiny\,(0.0053)} & +7.91{\tiny\,(0.0000)} & +8.34{\tiny\,(0.0000)} \\
    $\Delta$ w.r.t.\ RePO
    & +2.49{\tiny\,(0.0008)} & +4.44{\tiny\,(0.0000)} & +6.31{\tiny\,(0.0072)} & +5.58{\tiny\,(0.0234)} & +4.29{\tiny\,(0.0038)} & +3.83{\tiny\,(0.0001)} & +4.49{\tiny\,(0.0002)}
    & +9.84{\tiny\,(0.0000)} & +3.09{\tiny\,(0.0024)} & +24.65{\tiny\,(0.0000)} & +12.53{\tiny\,(0.0000)} \\
    $\Delta$ w.r.t.\ GRESO
    & +1.70{\tiny\,(0.0011)} & +5.57{\tiny\,(0.0006)} & +6.24{\tiny\,(0.0077)} & +6.24{\tiny\,(0.0052)} & +3.23{\tiny\,(0.0100)} & +3.72{\tiny\,(0.0002)} & +4.45{\tiny\,(0.0002)}
    & +0.32{\tiny\,(0.0032)} & +0.90{\tiny\,(0.1806)} & +1.92{\tiny\,(0.0045)} & +1.05{\tiny\,(0.0077)} \\
    $\Delta$ w.r.t.\ DAPO
    & +2.71{\tiny\,(0.0000)} & +2.37{\tiny\,(0.0011)} & +2.54{\tiny\,(0.0220)} & +8.00{\tiny\,(0.0010)} & +7.87{\tiny\,(0.0012)} & +4.89{\tiny\,(0.0001)} & +4.73{\tiny\,(0.0001)}
    & +3.84{\tiny\,(0.0013)} & +4.39{\tiny\,(0.0019)} & +3.68{\tiny\,(0.0001)} & +3.97{\tiny\,(0.0001)} \\

    \midrule
    \rowcolor{ModelRow}
    \multicolumn{12}{c}{\textbf{Qwen3-1.7B-Base} trained with \textbf{DAPO-Math-17k}} \\
    \midrule
    GRPO
    & 67.50{\tiny$\pm$0.79} & 25.29{\tiny$\pm$0.31} & 9.06{\tiny$\pm$0.63} & 4.48{\tiny$\pm$0.21} & 23.40{\tiny$\pm$1.39} & 29.57{\tiny$\pm$0.37} & 26.55{\tiny$\pm$0.26}
    & 33.62{\tiny$\pm$0.09} & 21.90{\tiny$\pm$1.04} & 12.71{\tiny$\pm$0.17} & 22.74{\tiny$\pm$0.32} \\
    RePO
    & 67.27{\tiny$\pm$0.55} & 23.78{\tiny$\pm$0.12} & 7.60{\tiny$\pm$1.50} & 5.94{\tiny$\pm$0.71} & 22.57{\tiny$\pm$1.39} & 29.00{\tiny$\pm$0.43} & 26.03{\tiny$\pm$0.29}
    & 36.58{\tiny$\pm$0.73} & 25.25{\tiny$\pm$0.99} & \textbf{30.37}{\tiny$\pm$0.41} & 30.73{\tiny$\pm$0.32} \\
    GRESO
    & 67.60{\tiny$\pm$0.05} & 23.76{\tiny$\pm$0.20} & 8.19{\tiny$\pm$1.37} & 6.39{\tiny$\pm$0.79} & 23.61{\tiny$\pm$1.96} & 30.91{\tiny$\pm$0.15} & 26.74{\tiny$\pm$0.09}
    & 35.68{\tiny$\pm$0.35} & 27.15{\tiny$\pm$1.24} & 26.67{\tiny$\pm$0.18} & 29.83{\tiny$\pm$0.36} \\
    DAPO
    & 66.69{\tiny$\pm$0.13} & 24.34{\tiny$\pm$0.23} & 9.72{\tiny$\pm$0.86} & 5.42{\tiny$\pm$0.90} & 19.72{\tiny$\pm$0.45} & 28.70{\tiny$\pm$0.06} & 25.77{\tiny$\pm$0.33}
    & 33.71{\tiny$\pm$0.69} & 24.49{\tiny$\pm$0.86} & 21.05{\tiny$\pm$0.12} & 26.42{\tiny$\pm$0.42} \\
    \midrule
    GRPO-DBB
    & \textbf{71.67}{\tiny$\pm$0.22} & \textbf{26.43}{\tiny$\pm$0.67} & \textbf{13.75}{\tiny$\pm$1.13} & \textbf{7.40}{\tiny$\pm$1.04} & \textbf{27.31}{\tiny$\pm$1.90} & \textbf{33.36}{\tiny$\pm$0.50} & \textbf{29.98}{\tiny$\pm$0.28}
    & \textbf{40.81}{\tiny$\pm$0.11} & \textbf{27.89}{\tiny$\pm$1.46} & 29.67{\tiny$\pm$0.89} & \textbf{32.79}{\tiny$\pm$0.79} \\
    $\Delta$ w.r.t.\ GRPO
    & +4.17{\tiny\,(0.0004)} & +1.14{\tiny\,(0.0419)} & +4.69{\tiny\,(0.0063)} & +2.92{\tiny\,(0.0037)} & +3.90{\tiny\,(0.0158)} & +3.79{\tiny\,(0.0007)} & +3.43{\tiny\,(0.0001)}
    & +7.19{\tiny\,(0.0000)} & +5.98{\tiny\,(0.0085)} & +16.97{\tiny\,(0.0000)} & +10.05{\tiny\,(0.0002)} \\
    $\Delta$ w.r.t.\ RePO
    & +4.40{\tiny\,(0.0002)} & +2.64{\tiny\,(0.0015)} & +6.14{\tiny\,(0.0007)} & +1.46{\tiny\,(0.0219)} & +4.74{\tiny\,(0.0202)} & +4.35{\tiny\,(0.0004)} & +3.95{\tiny\,(0.0002)}
    & +4.24{\tiny\,(0.0007)} & +2.64{\tiny\,(0.0110)} & $-$0.70{\tiny\,(0.8273)} & +2.06{\tiny\,(0.0007)} \\
    $\Delta$ w.r.t.\ GRESO
    & +4.07{\tiny\,(0.0000)} & +2.67{\tiny\,(0.0039)} & +5.56{\tiny\,(0.0001)} & +1.01{\tiny\,(0.0184)} & +3.69{\tiny\,(0.0209)} & +2.44{\tiny\,(0.0010)} & +3.24{\tiny\,(0.0001)}
    & +5.13{\tiny\,(0.0001)} & +0.74{\tiny\,(0.0354)} & +3.01{\tiny\,(0.0038)} & +2.96{\tiny\,(0.0001)} \\
    $\Delta$ w.r.t.\ DAPO
    & +4.98{\tiny\,(0.0000)} & +2.08{\tiny\,(0.0031)} & +4.03{\tiny\,(0.0119)} & +1.98{\tiny\,(0.0387)} & +7.58{\tiny\,(0.0018)} & +4.66{\tiny\,(0.0001)} & +4.22{\tiny\,(0.0002)}
    & +7.11{\tiny\,(0.0002)} & +3.39{\tiny\,(0.0119)} & +8.62{\tiny\,(0.0001)} & +6.38{\tiny\,(0.0004)} \\

    \bottomrule
    \end{tabular}
}
\end{table*}

Table~\ref{tab:main_results} reports Acc@8 results across six mathematical reasoning in-distribution benchmarks and three general reasoning out-of-distribution (OOD) benchmarks for five different RLVR algorithms.
For GRPO-DBB, the discounting factor $\lambda$ is set to $0.5$ for Qwen3-1.7B-Base and $0.75$ for Qwen3-8B-Base; these values are empirically selected based on performance trends under varying $\lambda$.
Further details are provided in Section~\ref{sec:impact_lambda}.


For in-distribution evaluation, GRPO-DBB consistently outperforms all four baselines across all benchmarks and model scales. Compared to GRPO, Acc@8 improves by an average of 3.43 points for Qwen3-1.7B-Base and 2.32 points for Qwen3-8B-Base, and the improvements are statistically significant in nearly every cell ($p<0.05$).
This indicates that, in the RLVR setting, replacing GRPO's point estimator with the DBB estimator is more effective for in-distribution downstream tasks. When compared to history- or variance-collapse-aware baselines (RePO, GRESO, and DAPO), GRPO-DBB still achieves average ID Acc@8 gains of 3.24--4.49 (8B) and 2.96--4.22 (1.7B) points. Although RePO leverages diverse forms of historical information, including tokens, token probabilities, and rewards, while GRESO and DAPO modify the sampling/filtering strategy, GRPO-DBB relies solely on historical reward statistics. Despite this, its strong downstream performance suggests that the DBB estimator utilizes historical information both efficiently and effectively.

For out-of-distribution evaluation, GRPO-DBB consistently outperforms the baselines across nearly all OOD benchmarks and model scales. In particular, it achieves substantial average gains over GRPO of 10.05 and 8.34 points for Qwen3-1.7B-Base and Qwen3-8B-Base, respectively. The only non-significant cell is BBH against RePO on the 1.7B model, where the two methods are statistically indistinguishable. These results demonstrate that the DBB estimation provides a strong advantage in improving generalization, which can be attributed to the characteristic of the Bayesian approach to account for uncertainty.
